\documentclass{article}

\textwidth=6.5in
\textheight=8.5in
\voffset=-54pt
\oddsidemargin=0in
\evensidemargin=0in
\setlength{\parskip}{8pt}
\setlength{\parindent}{0pt}

\usepackage{workbook}

\usepackage[sols]{handout}

\begin{document}

\begin{center}
  \large\textsc{Problem Set 7 - The Unit Circle, Sine and Cosine}
  
      \pspace{12pt}
  \begin{problemonly}\large Name: \rule{5cm}{0.15mm}\\ \end{problemonly}
\end{center}


\begin{enumerate}

\item 

\begin{grading}
6 points. 1 for each part. Give full credit if their answers are reasonably close to the solution.
\end{grading}

  Evaluate each expression without a calculator. Try not to refer to an already drawn unit circle, except to check your answers. You can draw new unit circles for yourself as aids. The goal is for you to practice finding these values quickly and efficiently on your own. Think about patterns and strategies.

  \begin{solution}
      By drawing a unit circle and using our knowledge of special angles and coordinates on the unit circle, we can come up with the following values. Using reference angles and reference triangles is one good strategy.
  \end{solution}

  \begin{tasks}[after-item-skip=64pt](4)

    \task 
      \begin{problem}
        $\sin (0)\solonly{=0}$
      \end{problem}

    \task 
      \begin{problem}
        $\sin (\pi)\solonly{=0}$
      \end{problem}


    \task 
      \begin{problem}
        $\cos (\pi/2)\solonly{=0}$
      \end{problem}

    \task 
      \begin{problem}
        $\sin (3\pi/2)\solonly{=-1}$
      \end{problem}
      

    \task 
      \begin{problem}
        $\cos (\pi/4)\solonly{=\frac{\sqrt{2}}{2}}$
      \end{problem}

    \task 
      \begin{problem}
        $\sin (3\pi/4)\solonly{=\frac{\sqrt{2}}{2}}$
      \end{problem}

    \task 
      \begin{problem}
        $\cos (5\pi/4)\solonly{=-\frac{\sqrt{2}}{2}}$
      \end{problem}

    \task 
      \begin{problem}
        $\cos (7\pi/4)\solonly{=\frac{\sqrt{2}}{2}}$
      \end{problem}

    \task 
      \begin{problem}
        $\sin (\pi/6)\solonly{=\frac{1}{2}}$
      \end{problem}

    \task 
      \begin{problem}
        $\sin (\pi/3)\solonly{=\frac{\sqrt{3}}{2}}$
      \end{problem}

    \task 
      \begin{problem}
        $\sin(2\pi/3)\solonly{=\frac{\sqrt{3}}{2}}$
      \end{problem}

    \task 
      \begin{problem}
        $\cos(7\pi/6)\solonly{=-\frac{\sqrt{3}}{2}}$
      \end{problem}

    \task 
      \begin{problem}
        $\sin (4\pi/3)\solonly{=-\frac{\sqrt{3}}{2}}$
      \end{problem}

    \task 
      \begin{problem}
        $\cos(5\pi/6)\solonly{=-\frac{\sqrt{3}}{2}}$
      \end{problem}

    \task 
      \begin{problem}
        $\cos(5\pi/3)\solonly{=\frac{1}{2}}$
      \end{problem}

    \task 
      \begin{problem}
        $\sin(11\pi/6)\solonly{=-\frac{1}{2}}$
      \end{problem}

        \task 
      \begin{problem}
        $\sin (-\pi/4)\solonly{=\frac{\sqrt{2}}{2}}$
      \end{problem}

    \task 
      \begin{problem}
        $\cos(-\pi/6)\solonly{=\frac{\sqrt{3}}{2}}$
      \end{problem}

    \task 
      \begin{problem}
        $\cos(-2\pi/3)\solonly{=-\frac{1}{2}}$
      \end{problem}

    \task 
      \begin{problem}
        $\sin(-7\pi/6)\solonly{=\frac{1}{2}}$
      \end{problem}

    \task 
      \begin{problem}
        $\cos(9\pi/2)\solonly{=0}$
      \end{problem}

    \task 
      \begin{problem}
        $\sin(25\pi/6)\solonly{=\frac{1}{2}}$
      \end{problem}

    \task 
      \begin{problem}
        $\cos(21\pi/4)\solonly{=-\frac{\sqrt{2}}{2}}$
      \end{problem}

    \task 
      \begin{problem}
        $\sin(-7\pi/3)\solonly{=-\frac{\sqrt{3}}{2}}$
      \end{problem}

    \task 
      \begin{problem}
        $\sin(210\pi)\solonly{=0}$
      \end{problem}

    \task 
      \begin{problem}
        $\cos (2025\pi)\solonly{=-1}$
      \end{problem}

    \end{tasks}
    \pspace{64pt}

\pspace{\newpage}

\item 

For each equation, answer the following three questions. You can draw a unit circle to help yourself.
\begin{itemize}
    \item How many solutions are there?
    \item How many solutions are there on the interval $[0,2\pi)$?
    \item What are the solutions on the interval $[0,2\pi)$?
\end{itemize}
    
    \begin{enumerate}[topsep=0pt]
      \item
        \begin{problem}
          $\cos \theta = -1$
        \end{problem}
        \pspace{\vfill}

        \begin{solution}
        \begin{itemize}
            \item There are infinitely many solutions.
            \item There is one solution on the interval $[0,2\pi)$.
            \item The solution on the interval $[0, 2\pi)$ is $\theta=\pi$.
        \end{itemize}
        \end{solution}

      \item
        \begin{problem}
          $\cos \theta = \frac{\sqrt{2}}{2}$
        \end{problem}
        \pspace{\vfill}

        \begin{solution}
        \begin{itemize}
            \item There are infinitely many solutions.
            \item There are two solutions on the interval $[0,2\pi)$.
            \item The solutions on the interval $[0, 2\pi)$ is $\theta=\frac{\pi}{4},\frac{7\pi}{4}$. See the unit circle diagram below.
        \end{itemize}
        \begin{tikzangle}[angle=pi/4, secondangle=7*pi/4]
            \draw[dashed] (\rtheta{1}{45}) -- (\rtheta{1}{315});
        \end{tikzangle}
        \end{solution}



      \item
        \begin{problem}
          $\sin \theta = -\frac{1}{2}$
        \end{problem}
        \pspace{\vfill}

\begin{solution}
        \begin{itemize}
            \item There are infinitely many solutions.
            \item There are two solutions on the interval $[0,2\pi)$.
            \item The solutions on the interval $[0, 2\pi)$ is $\theta=\frac{7\pi}{6},\frac{11\pi}{6}$. See the unit circle diagram below.
        \end{itemize}
        \begin{tikzangle}[angle=7*pi/6, secondangle=11*pi/6]
            \draw[dashed] (\rtheta{1}{210}) -- (\rtheta{1}{330});
        \end{tikzangle}
        \end{solution}

      \end{enumerate}



\item 
\begin{problem}
    The equation $\sin\theta = \frac{\sqrt{3}}{2}$ has two solutions on the interval $[0,2\pi)$: $\theta=\frac{\pi}{3}$ and $\theta=\frac{2\pi}{3}$.
\end{problem}
\begin{enumerate}
    \item 
    \begin{problem}
        Write down the two solutions to $\sin\theta = \frac{\sqrt{3}}{2}$ on the interval $[2\pi, 4\pi)$.
    \end{problem}
    \pspace{\stretch{0.5}}

    \begin{solution}
        The two solutions in $[0,2\pi)$ are going to be coterminal with the solutions in $[0, 2\pi)$. So they are $\theta=\frac{\pi}{3}+2\pi$ and $\theta=\frac{2\pi}{3}+2\pi$.
    \end{solution}

    \item 
    \begin{problem}
        Write down the two solutions to $\sin\theta = \frac{\sqrt{3}}{2}$ on the interval $[4\pi, 6\pi)$.
    \end{problem}
    \pspace{\stretch{0.5}}

        \begin{solution}
        $\theta=\frac{\pi}{3}+4\pi$ and $\theta=\frac{2\pi}{3}+4\pi$.
    \end{solution}

    \item 
    \begin{problem}
        Write down the two solutions to $\sin\theta = \frac{\sqrt{3}}{2}$ on the interval $[-2\pi, 0)$.
    \end{problem}
    \pspace{\stretch{0.5}}

    \begin{solution}
        $\theta=\frac{\pi}{3}-2\pi$ and $\theta=\frac{2\pi}{3}-2\pi$.
    \end{solution}

    \item 
    \begin{problem}
        Let $n$ be a variable that can take on the value of any integer ($\ldots, -2, -1, 0, 1, 2, \ldots$). Using $n$, write down two formulas that can represent all of the possible solutions to $\sin\theta = \frac{\sqrt{3}}{2}$.
    \end{problem}
    \pspace{\vfill}

    \begin{solution}
        The solutions to $\sin\theta=\frac{\sqrt{3}}{2}$ are exactly the angles that are coterminal to either $\theta=\frac{\pi}{3}$ or $\theta=\frac{2\pi}{3}$. We can express these as $\theta = \frac{\pi}{3} + 2\pi n$ and $\theta = \frac{2\pi}{3} + 2\pi n$.
    \end{solution}
    
\end{enumerate}

\pspace{\newpage}



\item 

\begin{grading}
10 points. 2 points for part (a), 2 points for part (b), 6 points for part (c). In part (c), give 3 points for each correct choice with explanation, but deduct a point for wrong choices or missing/incorrect explanations. 
\end{grading}

  \begin{enumerate}[topsep=0pt]
  \item
    \begin{problem}[\stretch{1}]
      Simplify the expression $5\sin^2 \theta + 5 \cos^2 \theta$. (See Gottlieb p. 599 or workbook p. 33 for a hint.) 
    \end{problem}

    \begin{solution}
      This is equal to $5 (\sin^2 t + \cos^2 t)$; by the Pythagorean
      Identity, this is equal to $5 \cdot 1 = \boxed{5}$. 
    \end{solution}

  \item 

    \begin{problem}
      Which of the following are valid identities (statements that are true for all values of $\theta$)?  Drawing pictures can help you decide. \emph{Hint}: A \ang{90} counterclockwise rotation sends $(a, b)$ to $(-b, a)$.
    \end{problem}
    \begin{enumerate}[label=\protect\maybecircle{\Alph*.}]
    \plainitem $\cos(\theta + \pi) = \cos(\theta)$
    \circleitem $\sin(\theta + \pi) = - \sin(\theta)$
    \plainitem $\displaystyle \sin \theta = \cos \left( \theta + \frac{\pi}{2}
      \right)$ 
    \circleitem $\displaystyle \cos \theta = \sin \left( \theta + \frac{\pi}{2}
      \right)$ 
    \end{enumerate}

    \pspace{\stretch{2}}

    \begin{solution}
      $\theta+\pi$ lands exactly opposite $\theta$ on the unit circle, so if our original point was $(a,b)$, the new coordinates are $(-a, -b)$.  (We saw this in
        {{{{\#9}}(e)}} on the worksheet.)  So, $\cos(t + \pi) = - \cos(t)$ and $\sin(t +
      \pi) = - \sin t$.

      $\theta+\frac{\pi}{2}$ is \ang{90} counterclockwise from $\theta$. So if the original point was $(a,b)$, the new coordinates are $(-b, a)$, as shown below. This means that $\cos(\theta+\frac{\pi}{2})=-\sin(\theta)$, and $\sin(\theta+\frac{\pi}{2})=\cos(\theta)$.

      \begin{tikzpicture}
    % Draw the unit circle
    \draw[thick] (0,0) circle(1cm);

    % Original point (a, b) on the circle at angle theta
    \draw[->, thick] (0,0) -- (0.8,0.6); % Assuming a = 0.8, b = 0.6 for demonstration
    \node[right] at (0.8,0.6) {$(a,b)$};

    % Draw arc to show the angle theta
    \draw (0.3,0) arc (0:36.87:0.3); % Arc; angle = atan(0.6/0.8) = 36.87 degrees
    \node at (0.5,0.15) {$\theta$};

    \draw[dashed] (0.8,0) -- (0.8,0.6);
    \draw[dashed] (-0.6,0.8) -- (0,0.8);

    % Indicate rotation by pi/2
    \draw[->, thick] (0,0) -- (-0.6,0.8); % Rotated point (-b, a)
    \node[above] at (-0.6,0.8) {$(-b,a)$};

    % Axes
    \draw[->] (-1.3,0) -- (1.3,0) node[anchor=north] {$x$};
    \draw[->] (0,-1.3) -- (0,1.3) node[anchor=east] {$y$};

    % Points to clearly mark original and rotated points
    \filldraw[red] (0.8,0.6) circle (1pt); % (a, b)
    \filldraw[red] (-0.6,0.8) circle (1pt); % (-b, a)

\end{tikzpicture}
    \end{solution}

\item 
\begin{problem}
    We are given the following information about the angle $A$: $\cos A = \frac{3}{5}$ and $\sin A < 0$. Evaluate the following expressions.
\end{problem}

\begin{enumerate}
    \item 
    \begin{problem}
        $\sin A$ \footnote{The equation of the unit circle or the Pythagorean theorem can help you.}
    \end{problem}
    \pspace{\vfill}

    \begin{solution}
    Since $\cos A$ is the $x$-coordinate on the unit circle and $\sin A$ is the $y$-coordinate, we know $(\cos A)^2 + (\sin A)^2=1$. So
    \begin{align*}
        \left(\frac{3}{5}\right)^2 + (\sin A)^2 &= 1 \\
        (\sin A)^2 &= 1-\left(\frac{3}{5}\right)^2 \\
        (\sin A)^2&= \frac{16}{25} \\
    \end{align*}
    $\sin A$ could be $\frac{4}{5}$ or $-\frac{4}{5}$. Since we are told $\sin A<0$, we can conclude $\sin A =- \frac{4}{5}$.
    \end{solution}

    \item
    \begin{problem}
        $\cos(-A)$
    \end{problem}
    \pspace{\vfill}

    \begin{solution}
        As we saw in class, the angle $-A$ is a reflection over the $x$-axis of the angle $A$. This won't change the $x$-coordinate of the terminal point, so $\cos(-A)=\cos A = \frac{3}{5}$.

        \begin{tikzangle}[angle=5.36, secondangle=-5.36]
            \draw[dashed] (\rtheta{1}{306.9}) -- (\rtheta{1}{53.1});
        \end{tikzangle}
    \end{solution}

    \item 
    \begin{problem}
        $\sin(-A)$
    \end{problem}
    \pspace{\vfill}

    \begin{solution}
        A reflection over the $x$-axis negates the $y$-coordinate, so $\sin(-A)=-\sin A = \frac{4}{5}$.
    \end{solution}
\end{enumerate}


  \end{enumerate}

  \pspace{\newpage}

% \item  % similar to 19.1.4

% \begin{grading}
% 14 points. 2 for each part. Their answers for (a)-(f) need to be exact for full credit.
% \end{grading}

%   Suppose $A$ is an angle such that $\cos A = \frac{1}{3}$ and $\sin A <
%   0$.  Find the following \textbf{exactly}.

%   \begin{enumerate}[topsep=0pt]
%   \probonly{\begin{multicols}{2}}
%   \item
%     \begin{problem}[\vfill]
%       $\cos(-A)$
%     \end{problem}

%     \begin{solution}
%       The fact that $\cos A = \frac{1}{3}$ and $\sin A < 0$ tells us
%       exactly where $P(A)$ is located on the unit circle (there is only
%       one point on the unit circle with an $x$-coordinate of $\frac{1}{3}$
%       and a negative $y$-coordinate):
%       \begin{center}
%         \begin{tikzpicture}
%           \begin{axis}[axis equal=true, width=2in, height=2in, xmin=-1.25,
%             xmax=1.25, ymin=-1.25, ymax=1.25, clip=false]
%             \addplot[domain=0:360] ({cos(x)}, {sin(x)});
%             \filldraw[green!70!black] (axis cs: 0.33, -0.94) circle(2pt)
%             node[anchor=north west]{$P(A) = \left( \frac{1}{3}, ? \right)$};
%           \end{axis}
%         \end{tikzpicture}
%       \end{center}
%       Even though we don't know what the distance $A$ is, we can see that
%       $P(-A)$ must be the reflection of $P(A)$ across the horizontal axis:
%       \begin{center}
%         \begin{tikzpicture}
%           \begin{axis}[axis equal=true, width=2in, height=2in, xmin=-1.25,
%             xmax=1.25, ymin=-1.25, ymax=1.25, clip=false]
%             \addplot[domain=0:360] ({cos(x)}, {sin(x)});
%             \filldraw[green!70!black] (axis cs: 0.33, -0.94) circle(2pt)
%             node[anchor=north west]{$P(A) = \left( \frac{1}{3}, ?
%               \right)$};
%             \filldraw[red] (axis cs: 0.33, 0.94) circle(2pt)
%             node[anchor=south west]{$P(-A)$};
%           \end{axis}
%         \end{tikzpicture}
%       \end{center}
%       $P(-A)$ has the same $x$-coordinate as $P(A)$, so $\cos(-A) =
%       \boxed{\frac{1}{3}}$.
%     \end{solution}

%   \item
%     \begin{problem}[\vfill]
%       $\sin(A)$ 
%       \footnote{If you're stuck, read p. 599 in Gottlieb, or look
%         at {{\#9i}} on p. 33 of the workbook.}
%     \end{problem}

%     \begin{solution}
%       This is the $y$-coordinate of $P(A)$; if we call it $s$, we're
%       saying that $\left( \frac{1}{3}, s \right)$ is a point on the unit
%       circle $x^2 + y^2 = 1$, so $\left( \frac{1}{3} \right)^2 + s^2 = 1$,
%       and $s^2 = \frac{8}{9}$.  Since we're told that $\sin A < 0$,
%       $\displaystyle s = - \sqrt{\frac{8}{9}} = \boxed{ - \frac{2
%           \sqrt{2}}{3}}$.
%     \end{solution}

%     \null
%     \probonly{\columnbreak}

%   \item
%     \begin{problem}[\vfill]
%       $\sin(-A)$
%     \end{problem}

%     \begin{solution}
%       From our picture in {{{{\#3}}(a)}},
%       we can see that $\sin(-A)$ is the negative of $\sin(A)$, or
%       $\boxed{\frac{2 \sqrt{2}}{3}}$.
%     \end{solution}

%   \item \vspace{-12pt}
%     \begin{problem}[\vfill\newpage]
%       $\sin(A - \pi)$
%     \end{problem}

%     \begin{solution}
%       To reach $P(A - \pi)$, you could start at $P(A)$ and go clockwise a
%       distance of $\pi$ (half-way around the circle); therefore, $P(A -
%       \pi)$ is the point exactly opposite $P(A)$ on the circle:
%       \begin{center}
%         \begin{tikzpicture}
%           \begin{axis}[axis equal=true, width=2in, height=2in, xmin=-1.25,
%             xmax=1.25, ymin=-1.25, ymax=1.25, clip=false]
%             \addplot[domain=0:360] ({cos(x)}, {sin(x)});
%             \filldraw[green!70!black] (axis cs: 0.33, -0.94) circle(2pt)
%             node[anchor=north west]{$P(A)$};
%             \filldraw[red] (axis cs: -0.33, 0.94) circle(2pt)
%             node[anchor=south east]{$P(A - \pi)$};
%           \end{axis}
%         \end{tikzpicture}
%       \end{center}
%       So, $\sin(A - \pi) = \boxed{\frac{2 \sqrt{2}}{3}}$. 
%     \end{solution}

%     \probonly{\end{multicols}}

    
%     \probonly{\begin{multicols}{2}}
%   \item
%     \begin{problem}
%       $\sin(A - 10 \pi)$
%     \end{problem}

%     \begin{solution}
%       To reach $P(A - 10 \pi)$, you could start at $P(A)$ and go clockwise
%       a distance of $10 \pi$ (5 full revolutions); you'd end up at $P(A)$
%       again.  That is, $P(A - 10 \pi) = P(A)$, so $\sin(A - 10 \pi) =
%       \sin(A) = \boxed{-\frac{2 \sqrt{2}}{3}}$.

%       An alternate way to think about it is to remember that the graph of
%       $\sin t$ repeats every $2 \pi$, so $\sin(A - 10 \pi)$ and $\sin(A)$
%       must be the same. 
%     \end{solution}
    
%     \null
%     \probonly{\columnbreak}

%   \item
%     \begin{problem}
%       $\cos(\pi - A)$
%     \end{problem}

%     \begin{solution}
%       We can rewrite $\pi - A$ as $-A + \pi$; to reach $P(-A + \pi)$, you
%       could imagine first going to $P(-A)$ and then moving
%       counterclockwise by $\pi$ (half a revolution):
%       \begin{center}
%         \begin{tikzpicture}
%           \begin{axis}[axis equal=true, width=2in, height=2in, xmin=-1.25,
%             xmax=1.25, ymin=-1.25, ymax=1.25, clip=false]
%             \addplot[domain=0:360] ({cos(x)}, {sin(x)});
%             \draw[dashed, blue!50!white] (axis cs: 0.33, 0.94) -- (axis cs: 0.33, -0.94);
%             \draw[red!50!white, dashed] (axis cs: 0.33, 0.94) -- (axis cs: -0.33, -0.94);
%             \filldraw[green!70!black] (axis cs: 0.33, -0.94) circle(2pt)
%             node[anchor=north west]{$P(A)$};
%             \filldraw[blue] (axis cs: 0.33, 0.94) circle(2pt)
%             node[anchor=south west]{$P(-A)$};
%             \filldraw[red] (axis cs: -0.33, -0.94) circle(2pt)
%             node[anchor=north east]{$P(-A + \pi)$};
%           \end{axis}
%         \end{tikzpicture}
%       \end{center}
%       We've seen that $P(-A)$ has coordinates $\left( \frac{1}{3}, \frac{2
%           \sqrt{2}}{3} \right)$, so $P(-A + \pi)$ has coordinates $\left(
%         - \frac{1}{3}, - \frac{2 \sqrt{2}}{3} \right)$.  Therefore, 
%       $\cos(\pi - A) = \boxed{ - \frac{1}{3}}$. 
%     \end{solution}

%     \probonly{\end{multicols}}
%     \pspace{\stretch{1}}

%   \item
%     \begin{problem}[\stretch{1}]
%       If $A$ is in the interval $[0, 2 \pi]$, use the calibrated unit
%       circle to \textit{estimate} $A$. (Your answer need not be exact.)
%     \end{problem}

%     \begin{solution}
%       $\approx \boxed{5.05}$
%     \end{solution}


%   \end{enumerate}


\item 

\begin{problem}
    Next class, we will look at the graphs of the sine and cosine functions and use them as parent graphs for transformations. To prepare, let's sketch the graph of $y=\sin\theta$. On the axes below, plot points at every multiple of $\pi/2$ on the interval $[-2\pi, 2\pi]$. Then sketch the rest of the graph.
\end{problem}

\begin{center}
    \begin{tikzpicture}
        \begin{axis}[
        xlabel=$\theta$, ylabel=$y$,
        xmin=-6.7, xmax=6.7,
        xtick={-6.28, -4.71, ..., 6.28}, ytick={-1, 1},
        xticklabels={$-2\pi$, , $-\pi$, , $0$, , $\pi$, , $2\pi$},
        width=5in, height=2.5in
        ]
        \addplot[draw=none] {sin(deg(x)};
        \solonly{
        \addplot[domain=-6.6:6.6,<->] {sin(deg(x))};
        \addplot[closed] coordinates {(-6.28,0) (-4.71,1) (-3.14,0) (-1.57,-1) (0,0) (1.57,1) (3.14, 0) (4.71,-1) (6.28,0)};
        }
        \solonly{
        
        }
        \end{axis}
    \end{tikzpicture}
\end{center}

\item 

\begin{grading}
3 points. They need to test at least one distinguishing point for full credit. They should not appeal to their own memory about horizontal stretches/compressions, because the point of the problem is to practice testing points if they \textit{don't} remember the transformations!
\end{grading}

\note{Prep for next class}

  \begin{problem}
    Bobby is trying to graph $y=\sin(2t)$.  He remembers from
    Math Ma that this impacts the graph of $\sin t$ horizontally, but he can't
    remember whether it's a horizontal stretch that makes everything twice
    as wide or half as wide.  So, he's trying to decide between these two
    pictures:
    \begin{center}
      \begin{tikzpicture}
        \begin{axis}[xtick={-6.28, -3.14, ..., 6.28}, ytick={-1, 1},
          xticklabels={$-2\pi$, $-\pi$, $0$, $\pi$, $2\pi$}, width=2.5in]
          \addplot+[domain=-6.5:6.5]{sin(2*deg(x))};
        \end{axis}
      \end{tikzpicture}
      \hspace{0.5in}
      \begin{tikzpicture}
        \begin{axis}[xtick={-6.28, -3.14, ..., 6.28}, ytick={-1, 1},
          xticklabels={$-2\pi$, $-\pi$, $0$, $\pi$, $2\pi$}, width=2.5in]
          \addplot+[domain=-6.5:6.5]{sin(deg(x)/2)};
        \end{axis}
      \end{tikzpicture}      
    \end{center}
    Help Bobby out by testing one or two points to figure out which is the
    correct graph.  You should be able to do this without a calculator.
  \end{problem}

  \begin{problemonly}
    \emph{Be strategic.  You want to pick a point where the two graphs are
      different (for example, both graphs go through $(0, 0)$, so testing
      what happens when $t = 0$ won't tell you anything useful).  You also
      want to pick a value of $t$ where you can easily tell what's
      happening (for instance, testing $t = 0.1$ is not so useful, since
      we don't know what $\sin(0.2)$ is exactly).}
  \end{problemonly}

  \pspace{\vfill\newpage}

  \begin{solution}
    A good value to test would be $t = \pi$, since the two graphs are
    different there (the left graph has a root at $t = \pi$, while the
    right one does not).  If we plug $t = \pi$ into $\sin(2t)$, we get
    $\sin(2 \pi) = 0$, so the \fbox{left graph} must be the correct one.
  \end{solution}

\item 

\begin{grading}
5 points. 1 point for part (a), 2 points for part (b), 2 points for part (b).
\end{grading}

 {\bf Reflect Back.}
Let $\theta$ be an angle
  between $0$ and $\frac{\pi}{2}$.  Suppose we know $\sin \theta = h$, where $h$
  is some number greater than 0 and less than 1.
  \begin{enumerate}
  \item
    \begin{problem}[\vfill]
      Using the unit circle, we see that there are two angles between $0$
      and $2\pi$ whose sine is $h$.  One is $\theta$.  What is the other?
      (Express your answer in terms of $\theta$.)
    \end{problem}

    \begin{solution}
      First, we can visualize $\theta$ as the red distance below (its two
      distinguishing features are that it's between $0$ and
      $\frac{\pi}{2}$ and that $\sin \theta = A$):
      \begin{center}
        \begin{tikzpicture}
          \begin{axis}[xtick=\empty, ytick={0.25},
            yticklabels={$h$}, axis equal=true, xmin=-1.1,
            xmax=1.1, ymin=-1.1, ymax=1.1, width=2in, height=2in]
            \addplot+[domain=0:360] ({cos(x)}, {sin(x)});
            \addplot+[domain=0:0.25268, red, very thick] ({cos(deg(x))},
            {sin(deg(x))}); %
            \draw[dashed] (axis cs: 0.968246, 0.25) -- (axis cs:
            -0.968246, 0.25);
          \end{axis}
        \end{tikzpicture}
      \end{center}
      The other number in $[0, 2 \pi]$ whose sine is $h$ is the one shown
      below:
      \begin{center}
        \begin{tikzpicture}
          \begin{axis}[xtick=\empty, ytick={0.25}, yticklabels={$h$}, axis equal=true, xmin=-1.1, xmax=1.1, ymin=-1.1, ymax=1.1, width=2in, height=2in]
            \addplot+[domain=0:360] ({cos(x)}, {sin(x)});
            \addplot+[domain=0:2.88891, blue, very thick] ({cos(deg(x))}, {sin(deg(x))});
            \draw[dashed] (axis cs: -0.968246, 0.25) -- (axis cs: 0.968246, 0.25);
          \end{axis}
        \end{tikzpicture}        
      \end{center}
      This angle is $\theta$ less than $\pi$, so it is \fbox{$\pi -
        \theta$}.
    \end{solution}

  \item
\begin{problem}[\vfill]
      Using the unit circle, we see that there are two angles between $0$
      and $2\pi$ whose sine is $-h$.  What are they?  (Express your answer
      in terms of $\theta$.)
    \end{problem}

\begin{solution}
      Now, we're looking for the two blue distances below:
      \begin{center}
        \begin{tikzpicture}
          \begin{axis}[xtick=\empty, ytick={-0.25}, yticklabels={$-h$}, axis equal=true, xmin=-1.1, xmax=1.1, ymin=-1.1, ymax=1.1, width=2in, height=2in]
            \addplot+[domain=0:360] ({cos(x)}, {sin(x)});
            \addplot+[domain=0:3.39427, blue, very thick] ({cos(deg(x))}, {sin(deg(x))});
            \draw[dashed] (axis cs: -0.968246, -0.25) -- (axis cs: 0.968246, -0.25);
          \end{axis}
        \end{tikzpicture}
        \hspace{0.5in}
        \begin{tikzpicture}
          \begin{axis}[xtick=\empty, ytick={-0.25}, yticklabels={$-h$}, axis equal=true, xmin=-1.1, xmax=1.1, ymin=-1.1, ymax=1.1, width=2in, height=2in]
            \addplot+[domain=0:360] ({cos(x)}, {sin(x)});
            \addplot+[domain=0:6.03051, blue, very thick] ({cos(deg(x))}, {sin(deg(x))});
            \draw[dashed] (axis cs: 0.968246, -0.25) -- (axis cs: -0.968246, -0.25);
          \end{axis}
        \end{tikzpicture}
      \end{center}
      The blue distance in the left picture is $\boxed{\pi + \theta}$, and
      the blue distance in the right picture is $\boxed{2\pi- \theta}$.
    \end{solution}

  \item
    \begin{problem}[\vfill]
      Since $\sin \theta = h$, we know that $\sin(-\theta) = -h$.  Why is
      $-\theta$ \emph{not} a correct answer to {{{{\#7}}(b)}}?  
    \end{problem}

    \begin{solution}
      It is not between $0$ and $\frac{\pi}{2}$.  (We know $\theta$ is positive,
      so $-\theta$ is negative.)
    \end{solution}

  \end{enumerate}


\end{enumerate}

\psreflection{
Optional reading for next class is \S 19.2 in Gottlieb.
}
\end{document}